\scanpage{9}

Si $E$ est un \deleted{motif} \inserted{$\mathbb{Q}_{\ell}$-faisceau} sur $\unsure{\operatorname{Sp}|\operatorname{Spec}}\mathbb{F}_p$, \deleted{\unsure{donc}} on trouve donc
\[
L_{\check E}(t)=(-t)^{-n}\det f_E\,L_E(t^{-1})
\]
où $n=\operatorname{rang}E$ et s'il s'agit d'un \uline{motif} \inserted{de poids $\rho$} satisfaisant aux conj. de Weil
\[
\left\{\det f_E=\deleted{\varepsilon}\inserted{\varepsilon(E)}p^{\frac{n\rho}{2}}\right.
\]
avec $\varepsilon=\pm1$ [dépendant du motif $\bigwedge^n E$, \unsure{suivant|savoir} si $f$ y opère \unsure{positivement} ou non] ; noter que $\frac{n\rho}{2}$ est tjrs entier (i.e. $\rho$ impair $\Longrightarrow n$ pair).

On trouve donc
\[
L_{\check E}(t)=\varepsilon(E)(-p^{\rho/2}/t)^n L_E(t^{-1}).
\]

D'autre part, par Lefschetz, on aura
\[
\check E\simeq E(\rho)
\]
d'où
\[
\begin{aligned}
L_{\check E}(t)&=\left(L_E(t)*L_{\mathbb{Q}_{\ell}(\rho)}(t)\right)^{-1}\\
&=\left(L_E(t)*\frac{1}{1-p^{-\rho}t}\right)^{-1}=L_E(t)*(1-p^{-\rho}t)
\end{aligned}
\]
\[
L_{\check E}(t)=L_E\left(\frac{t}{p^{\rho}}\right)
\]

\mspage{8}
